\documentclass[12pt,a4paper]{article}
\usepackage{fontspec}
\IfFontExistsTF{Times New Roman}{\setmainfont{Times New Roman}}{\setmainfont{DejaVu Serif}[ItalicFont={DejaVu Serif},ItalicFeatures={FakeSlant=0.15},BoldItalicFont={DejaVu Serif Bold},BoldItalicFeatures={FakeSlant=0.15}]}
\setsansfont{DejaVu Sans}
\setmonofont[Scale=0.85]{DejaVu Sans Mono}
\usepackage[left=3cm,right=2cm,top=2cm,bottom=2cm]{geometry}
\usepackage{graphicx,booktabs,tabularx,array,longtable,float}
\usepackage{enumitem,setspace,titlesec,fancyhdr,xcolor}
\usepackage{listings,xurl}
\usepackage[hidelinks,unicode]{hyperref}
\usepackage{caption,tocloft}
\setlength{\cftsecnumwidth}{2.2em}
\setlength{\cftsubsecnumwidth}{2.7em}
\setstretch{1.15}
\setlength{\parindent}{1.1cm}
\setlength{\parskip}{3pt}
\setlength{\emergencystretch}{2em}
\setlist{nosep,leftmargin=1.6em,topsep=4pt}
\renewcommand{\contentsname}{Гарчиг}
\renewcommand{\figurename}{Зураг}
\renewcommand{\tablename}{Хүснэгт}
\captionsetup{font=small,labelfont=bf}
\titleformat{\section}{\large\bfseries}{\thesection.}{0.6em}{}
\titleformat{\subsection}{\normalsize\bfseries}{\thesubsection.}{0.6em}{}
\titlespacing*{\section}{0pt}{12pt}{7pt}
\titlespacing*{\subsection}{0pt}{8pt}{5pt}
\pagestyle{fancy}
\fancyhf{}
\fancyhead[L]{\small F.CSM302 Үйлдлийн систем}
\fancyhead[R]{\small Бие даалт №1}
\fancyfoot[C]{\thepage}
\setlength{\headheight}{16pt}
\renewcommand{\headrulewidth}{0.3pt}
\newcolumntype{Y}{>{\raggedright\arraybackslash}X}
\newcolumntype{P}[1]{>{\raggedright\arraybackslash}p{#1}}
\newcommand{\code}[1]{\texttt{#1}}
\newcommand{\docname}{Судалгаа, хэрэгжилт, туршилтын тайлан}
\lstset{basicstyle=\ttfamily\footnotesize,breaklines=true,columns=fullflexible,keepspaces=true,
 frame=single,rulecolor=\color{black!25},backgroundcolor=\color{black!3},showstringspaces=false,
 xleftmargin=4pt,xrightmargin=4pt,aboveskip=7pt,belowskip=7pt}
\hypersetup{pdfauthor={А.Тэмүүлэн B242270063},pdftitle={Linux Kernel Module for Listing Tasks}}

\renewcommand{\docname}{Хэрэглэгчийн болон operational гарын авлага}
\begin{document}
% User-provided MUST title-page layout, adapted for F.CSM302.
\begin{titlepage}
\begin{center}
{\large Шинжлэх Ухаан Технологийн Их Сургууль\par}
{Мэдээлэл, Холбооны Технологийн Сургууль\par}
\vspace{0.6cm}
\includegraphics[width=2.05cm]{figures/MUST_logo.png}\par
\vspace{0.8cm}
{\large F.CSM302\quad Үйлдлийн систем\par}
\vspace{1.25cm}
{\Large\bfseries Linux цөмийн модулиар процессын\par
жагсаалт ба холбоосыг дүрслэх\par}
\vspace{0.35cm}
{\normalsize Linux Kernel Module for Listing Tasks\par}
\vspace{0.65cm}
{\Large Бие даалт №1\par}
\vspace{0.25cm}
{\normalsize\docname\par}
\vspace{1.0cm}
\begin{minipage}{0.97\textwidth}
\setstretch{1.15}
\begin{tabularx}{\linewidth}{@{}P{4.0cm}Y@{}}
Шалгасан багш: & Г.Ганбат \\[0.2cm]
Багийн дугаар: & IV-2-3 \\[0.3cm]
Гүйцэтгэсэн оюутнууд: & А.Тэмүүлэн /B242270063/ \\[0.15cm]
& Ч.Балгандамба /B232270142/ \\[0.15cm]
& Д.Азмаа /B241940101/ \\[0.15cm]
& Н.Билэгт-Очир /B231910009/ \\[0.15cm]
& Т.Түвшинтөр /B190910820/ \\[0.15cm]
& Б.Өсөхбаяр /B210930901/ \\
\end{tabularx}
\end{minipage}
\vfill
{\large Улаанбаатар хот\par}
{\large 2026 он\par}
\end{center}
\end{titlepage}

\pagenumbering{arabic}

\section{Зорилго ба дэмжих орчин}
TaskGraph нь Linux-ийн process, thread-ийн холбоос, төлөвийг граф болон хүснэгтээр харуулна. Windows дээр жишээ эсвэл хадгалсан snapshot-оор UI-г туршиж болно. Бодит kernel module нь Linux дээр ажиллана.
\begin{table}[H]
\centering\small
\begin{tabularx}{\textwidth}{@{}P{3.2cm}Y@{}}\toprule
Горим & Шаардлагатай орчин \\\midrule
Windows UI & Python 3, Tcl/Tk; \code{run.bat}. \\
Linux UI & Python 3, \code{python3-tk}; \code{run.sh}. \\
Бодит module & Ubuntu 24.04 LTS, GCC/Make, ажиллаж буй kernel-ийн headers. Туршсан target: Linux 6.8.0-31. \\
Windows дахь Linux VM & VirtualBox; 2 CPU, 4 GB RAM, 40 GB диск санал болгоно. \\\bottomrule
\end{tabularx}
\end{table}

\section{Суулгах ба ажиллуулах}
\subsection{Windows дээр UI турших}
Python 3 (Tcl/Tk)-тай байхад \code{run.bat}-аар UI нээнэ. ZIP-ийг задалж эсвэл дараах командаар repository-г татна. Нэмэлт pip сан шаардахгүй.
\begin{lstlisting}
git clone https://github.com/Temka0511/FCSA_302-IV-2-3-LINUX-Kernel-Module-For-Task-Information.git taskgraph
cd taskgraph
run.bat
\end{lstlisting}
Модульгүй үед \textbf{SOURCE: FIXTURE} гэсэн бэлэн жишээ нээгдэнэ. Автомат тестийг \code{test.bat}-аар ажиллуулна. Windows дээр procfs-тэй холбоотой нэг тест алгасагдаж, бусад 22 тест ажиллана.

\subsection{Ubuntu дээр module ажиллуулах}
Ubuntu Terminal-д dependency суулгаад төслийн хавтас руу орно:
\begin{lstlisting}
sudo apt update
sudo apt install -y build-essential \
  linux-headers-$(uname -r) python3 python3-tk git
bash scripts/check_env.sh
make -C src/kernel
sudo insmod src/kernel/taskgraph.ko \
  max_tasks=4096 log_on_load=1
bash run.sh --kernel
\end{lstlisting}
UI-г энгийн хэрэглэгчээр ажиллуулна. Модуль ачаалах, буулгахад sudo хэрэгтэй. Дуусгаад \code{sudo rmmod taskgraph} командаар буулгана.
\textbf{Параметр:} \code{max\_tasks=16..16384}, анхны утга 4096; \code{log\_on\_load=1} нь LINEAR/DFS log хэвлэнэ. Өөр kernel ашиглавал өөрийн headers-тай дахин build хийнэ.

\clearpage
\section{UI-ийн удирдлага}
\begin{table}[H]
\centering\small
\begin{tabularx}{\textwidth}{@{}P{3.7cm}Y@{}}\toprule
Удирдлага & Үйлдэл \\\midrule
Эх үүсвэр & \code{kernel}: module; \code{procfs}: Linux оношилгоо; \code{file}: хадгалсан snapshot. \\
DFS / LINEAR & Task үзүүлэх дарааллыг солино. \\
Root PID & Тухайн task ба дэд мод; хоосон бол бүгд. \\
Threads & Нэмэлт thread-үүдийг харуулна эсвэл нууна. \\
Зангилааны хязгаар & 1--2000 хүртэл task харуулна. \\
Шинэчлэх / Авто шинэчлэлт & Эх үүсвэрийг дахин уншина; 500--10000 ms интервалтай. \\
Өмнөх / Дараах / Тоглуулах & Нэг snapshot-ийн traversal-ийг алхмаар үзүүлнэ. \\
Reset & Шүүлт, дараалал, playback-ийг эхний байдалд оруулна. \\
Файл нээх / Хадгалах & Бүтэн TSV snapshot импортлох, экспортлох. \\\bottomrule
\end{tabularx}
\end{table}
Зангилаа дээр дарвал PID, PPID, TGID, төлөв, дарааллын дэлгэрэнгүй гарна. Scroll болон дунд товчоор графыг шилжүүлнэ. Ногоон R, цэнхэр S, шар T/D, саарал I, ягаан Z/X төлөвийг илэрхийлнэ. Улбар хүрээ нь сонгосон task. Тасархай сум нь thread-ийн холбоос.

\section{Ашиглах жишээ}
\subsection{Жишээ 1: бэлэн модыг судлах}
\begin{enumerate}
\item \code{run.bat}-аар нээгээд Reset дарна.
\item DFS сонгож Тоглуулах дарна. Дараалал: 0, 1, 10, 11, 20, 21, 2.
\item Root PID=10 оруулаад Хэрэглэх дарна. 10, 11, 20, 21 гэсэн 4 task харагдана.
\item Threads-ийн чагтыг авахад 11 worker нуугдана.
\end{enumerate}
\subsection{Жишээ 2: kernel-ийн хадгалсан хэмжилт}
Файл нээхээр \code{docs/evidence/created.tsv} сонгож, Root PID=90 оруулна. 7 task бүхий мод гарна. Дараа нь \code{stopped.tsv} нээхэд root-ийн төлөв T болж шарласан байна. \code{after.tsv} нээхээс өмнө Root PID-г цэвэрлэнэ: workload аль хэдийн дууссан тул PID 90 байхгүй.
File горимын авто шинэчлэлт нь ижил файлыг дахин уншина; бодит процессын шинэ өөрчлөлтийг хэмжихгүй.

\clearpage
\subsection{Жишээ 3: бодит Linux процессын мод}
Module болон UI-г kernel горимд ажиллуулсны дараа өөр Terminal-д:
\begin{lstlisting}
export PYTHONPATH="$PWD/src"
python3 -m taskgraph.workload \
  --depth 2 --fanout 2 --seconds 45
\end{lstlisting}
ROOT\_PID мөрийн дугаарыг UI-д оруулж, Шинэчлэх, Хэрэглэх дарна. Энэ workload нь 1+2+4=7 process үүсгэнэ. Авто шинэчлэлтээр хугацаа дуусахад task-ууд алга болохыг ажиглана. Root өөрөө дуусвал талбарыг цэвэрлэнэ.
Depth=0..3, fanout=1..3, seconds=2..300. SIGSTOP/CONT турших бол зөвхөн өөрийн workload-ийн root PID-д дараах дохиог өгнө:
\begin{lstlisting}
kill -STOP ROOT_PID
kill -CONT ROOT_PID
\end{lstlisting}

\section{Operational шалгалт}
\subsection{Үндсэн build ба тест}
\begin{lstlisting}
make -C src/kernel
bash scripts/test.sh
sudo dmesg | grep 'taskgraph LINEAR' | head
sudo dmesg | grep 'taskgraph DFS' | head
ps -el
ps -eLf
cat /proc/taskgraph
\end{lstlisting}
Python тестийн хүлээгдсэн үр дүн нь Linux дээр 23 тест OK. LINEAR log-ийг process leader-уудтай, DFS log-ийг thread-үүдтэй харьцуулна. PID 0 нь \code{init\_task}. Хэмжилтийн агшин өөр учраас түр task, төлөвийн зөрүү гарч болно.
\subsection{Нэмэлт live ба QEMU тест}
\begin{lstlisting}
PYTHONPATH=src python3 scripts/verify_live.py
sudo apt install -y qemu-system-x86 busybox-static
PYTHONPATH=src python3 scripts/qemu_smoke.py \
  --kernel /boot/vmlinuz-$(uname -r)
\end{lstlisting}
Эхний скрипт ачаалсан module дээр 6 live шалгалт хийнэ. Хоёр дахь нь тохирох kernel image, build хийсэн module-оор QEMU дотор 12 шалгалт ажиллуулна. Image нь уншигдах боломжтой байна. Алдаа гарвал exit code 1 буцаана. Туршилтын log, JSON, TSV-г тусад нь хадгална.

\clearpage
\section{Түгээмэл алдаа ба хязгаарлалт}
\begin{table}[H]
\centering\small
\begin{tabularx}{\textwidth}{@{}P{4cm}Y@{}}\toprule
Алдаа & Шийдэл \\\midrule
\code{/proc/taskgraph} алга & Module build, \code{sudo insmod} хийсэн эсэхийг шалгана. \\
Invalid module format & \code{uname -r}-тэй таарсан headers суулгаж clean, build хийнэ. \\
Key was rejected & VM-ийн Secure Boot, module signing тохиргоог шалгана. \\
File exists & Хуучин module-ийг \code{sudo rmmod taskgraph}-аар буулгаад ачаална. \\
\code{tkinter} олдохгүй & Ubuntu-д \code{sudo apt install python3-tk}; Windows-д Python Tcl/Tk бүрэлдэхүүн суулгана. \\
DISPLAY алдаа & Desktop session-оос UI нээх эсвэл CLI ашиглана. \\
PID олдохгүй & Дууссан эсвэл тухайн snapshot-д байхгүй PID. Root талбарыг цэвэрлэнэ. \\
\code{truncated=1} & Capture cap хүрсэн. Module-ийг илүү их \code{max\_tasks}-аар дахин ачаална. \\
Графын task дутуу & View limit болон Threads шүүлтийг шалгана. \\\bottomrule
\end{tabularx}
\end{table}
Snapshot нь системийн бүх талбарыг нэг агшинд атомикаар хэмжихгүй. Task-ийн нэр нь богино \code{comm} утга. PID нь init namespace-ийнх тул module-ийг тусдаа Linux VM-д ашиглана. Windows-ийн fixture/file нь UI туршилтын горим юм.

\section{Төслийн файлууд}
\begin{table}[H]
\centering\small
\begin{tabularx}{\textwidth}{@{}P{4cm}Y@{}}\toprule
Зам & Агуулга \\\midrule
\code{src/kernel/} & C kernel module, Makefile. \\
\code{src/taskgraph/} & UI, parser, CLI, workload, launcher. \\
\code{tests/}, \code{examples/} & Автомат тест ба бэлэн өгөгдөл. \\
\code{scripts/} & Орчин, live болон QEMU шалгах скрипт. \\
\code{docs/evidence/} & Бодит kernel туршилтын log, JSON, TSV. \\
\code{docs/latex/} & Тайлан, Manual-ийн LaTeX эх ба ШУТИС-ийн лого. \\\bottomrule
\end{tabularx}
\end{table}
README-ийн Quick Start заавраар програмыг ажиллуулна. Эх код, тест, тайлан, гарын авлага, туршилтын нотолгоог GitHub-д хамт байршуулна.
\end{document}
